\documentclass[12pt]{article}

\usepackage{sbc-template}
\usepackage{graphicx,url}
\usepackage[utf8]{inputenc}
\usepackage[brazil]{babel}
\usepackage{amsmath}
\usepackage{booktabs}
\usepackage{xcolor}

% ============================================================
% PLACEHOLDERS — valores a preencher após rodar os experimentos.
% Cada placeholder \ph aparece em vermelho e negrito no PDF; basta localizar a macro.
% Fatos estruturais (tabela canônica, parâmetros, validador 21/21) são REAIS.
% Apenas os números que dependem da rodada estão como placeholder.
%
% Como gerar cada grupo:
%   baseline canônico .... py -m src.tools.report --nsga2 best_drift
%                          py -m src.tools.external_validation --canonical
%   AG escalar ........... py -m src.tools.report --evolved
%   fronteira/representantes  py -m src.tools.nsga2_plots ; pareto_metrics
%   estatística agregada . py -m src.tools.multi_run --algorithm nsga2 (e ga)
%   robustez ............. py -m src.tools.external_validation --nsga2 best_dominance
% ============================================================
\newcommand{\ph}[1]{\textcolor{red}{\textbf{[#1]}}}

\sloppy

\title{Balanceamento competitivo com preservação de identidade:\\
um estudo com Algoritmos Genéticos e NSGA-II em\\
arquétipos de jogos de luta}

\author{Matheus Gimenes de Souza\inst{1}, Prof. Dr. Everson Scherrer Borges\inst{1}}

\address{Instituto Federal do Espírito Santo (Ifes)\\
  Campus Cachoeiro de Itapemirim\\
  Cachoeiro de Itapemirim --- ES --- Brasil
  \email{contato.matheusgimenes@gmail.com, eversonborges@gmail.com}
}

\begin{document}

\maketitle

\begin{abstract}
  Balancing competitive games requires that distinct character archetypes remain
  viable without collapsing into a single dominant style. This work investigates
  whether a Genetic Algorithm can balance five fighting-game archetypes
  \emph{without destroying their functional identities}. Identity is
  \emph{measured}, not enforced: deviation from a canonical profile is a soft
  penalty, and the advantage cycle is reported \emph{post hoc}, which avoids a
  circular research question. Balance and identity are treated as two objectives,
  optimized both by a scalar GA and by NSGA-II. Aggregated over \ph{N} independent
  runs, the results show that the canonical roster is severely unbalanced
  (\ph{eq-can} balanced matchups); directly optimizing balance attains
  \ph{eq-ag} balanced matchups but homogenizes the roster, whereas a solution
  from the Pareto front reaches \ph{eq-bd} balanced matchups while \emph{increasing}
  the differentiation between archetypes. The front makes the trade-off between
  balance and identity explicit and measurable.
\end{abstract}

\begin{resumo}
  Balancear jogos competitivos exige que arquétipos distintos permaneçam viáveis
  sem colapsar em um único estilo dominante. Este trabalho investiga se um
  Algoritmo Genético é capaz de equilibrar cinco arquétipos de um jogo de luta
  \emph{sem destruir suas identidades funcionais}. A identidade é \emph{medida},
  não imposta: o desvio em relação a um perfil canônico é penalizado de forma
  suave, e o ciclo de vantagens é reportado \emph{post hoc}, o que evita uma
  pergunta de pesquisa circular. Equilíbrio e identidade são tratados como dois
  objetivos, otimizados tanto por um AG escalar quanto pelo NSGA-II. Agregados
  sobre \ph{N} execuções independentes, os resultados mostram que o roster
  canônico é fortemente desequilibrado (\ph{eq-can} confrontos equilibrados);
  otimizar o equilíbrio diretamente alcança \ph{eq-ag} confrontos equilibrados, mas
  homogeneíza o elenco, ao passo que uma solução da fronteira de Pareto atinge
  \ph{eq-bd} confrontos equilibrados \emph{aumentando} a diferenciação entre
  arquétipos. A fronteira torna explícito e mensurável o compromisso entre
  equilíbrio e identidade.
\end{resumo}

\section{Introdução}

A indústria de jogos digitais consolidou-se como um dos maiores mercados de
entretenimento do mundo, com receita global na casa das centenas de bilhões de
dólares e crescimento sustentado dos esportes eletrônicos
\cite{newzoo2025,pgb2026,hamari2017esports,reitman2020esports}. Em jogos
competitivos, a qualidade da experiência depende criticamente do
\emph{equilíbrio}: o grau em que as diferentes opções disponíveis ao jogador
---~personagens, classes, estratégias~--- oferecem chances comparáveis de vitória
\cite{adams2014fundamentals,schell2019art}. Um jogo desequilibrado tende a
degenerar em um \emph{metagame} estreito, no qual apenas algumas opções são
competitivamente viáveis, o que reduz a profundidade estratégica e o engajamento
\cite{sirlin2006playing,sweetser2005gameflow}.

O balanceamento, contudo, não se reduz a igualar taxas de vitória. Jogos
competitivos de qualidade são construídos sobre \emph{arquétipos}~--- identidades
funcionais fortes, como o agressor que sufoca o oponente de perto ou o
controlador que domina o espaço à distância~--- cuja distinção sustenta a
estrutura não transitiva (do tipo pedra-papel-tesoura) entre estilos
\cite{gingold2006rps,elias2012characteristics}. Decorre daí uma tensão central,
com frequência tratada de forma artesanal pelos projetistas: ao ajustar números
para igualar as taxas de vitória, corre-se o risco de \emph{homogeneizar} os
personagens, aproximando-os de um mesmo ponto e dissolvendo as identidades que
distinguem o jogo. Equilibrar é trivial quando todas as opções são iguais; o
desafio é equilibrar \emph{preservando a diferença}.

Este trabalho aborda essa tensão por meio de busca evolutiva. A pergunta de
pesquisa é:

\begin{quote}
\emph{Um Algoritmo Genético é capaz de atingir equilíbrio competitivo entre cinco
arquétipos distintos sem destruir suas identidades funcionais?}
\end{quote}

A decisão metodológica que sustenta a investigação é \textbf{não forçar a
identidade}. Os valores canônicos de cada arquétipo servem como semente da
população inicial e como \emph{baseline} de medição, nunca como restrição rígida;
o algoritmo evolui livremente. O desvio em relação ao perfil canônico é
\emph{penalizado de forma suave} (e exposto como objetivo de Pareto), enquanto o
\emph{ciclo de vantagens} canônico (quem vence quem) não é codificado em nenhuma
penalidade~--- é medido \emph{post hoc}. Essa escolha evita uma armadilha de
circularidade: se a preservação de identidade fosse imposta pela função de
avaliação, o algoritmo a ``preservaria'' apenas por ter sido recompensado para
tal, e o resultado não responderia à pergunta. Ao manter a identidade como algo
\emph{medido, não imposto}, sua preservação~--- ou sua ausência~--- torna-se um
achado genuíno. O princípio dialoga com o argumento de \cite{lehman2011novelty}
de que não codificar diretamente o objetivo pode revelar mais sobre o espaço de
busca do que otimizá-lo.

As contribuições deste artigo são: (i) uma formulação do balanceamento como
problema \emph{bi-objetivo} explícito~--- equilíbrio \emph{versus} preservação de
identidade~--- otimizado tanto por um AG escalar quanto pelo NSGA-II
\cite{deb2002nsga}; (ii) um conjunto de quatro instrumentos de medição
\emph{post hoc} da identidade~--- validação por invariantes de ranqueamento
(estruturais e comportamentais), assinatura comportamental, desvio numérico ao
canônico e diferenciação par a par~--- que quantificam a preservação de identidade
sem contaminar a função objetivo; e (iii) uma avaliação experimental que torna o
\emph{trade-off} mensurável e evidencia ser possível obter ganho substancial de
equilíbrio sem homogeneizar o elenco. As seções seguintes apresentam a
fundamentação teórica, a metodologia e os resultados.

\section{Fundamentação Teórica e Trabalhos Relacionados}

\subsection{Algoritmos Genéticos e otimização multiobjetivo}

Algoritmos Genéticos (AGs) são metaheurísticas de busca inspiradas na evolução
natural, nas quais uma população de soluções candidatas é iterativamente
submetida a seleção, recombinação (\emph{crossover}) e mutação, evoluindo em
direção a regiões de maior aptidão (\emph{fitness})
\cite{holland1975adaptation,goldberg1989genetic,mitchell1998introduction}. São
particularmente adequados a espaços de busca grandes, não lineares e sem gradiente
analítico~--- exatamente o caso de um espaço de parâmetros de \emph{game design}
avaliado por simulação. O compromisso entre exploração (diversificar) e explotação
(intensificar) é determinado por operadores e parâmetros como pressão de seleção,
taxa de mutação e elitismo
\cite{whitley1994tutorial,eiben1998exploration,eiben2015introduction}.

Quando o problema envolve objetivos conflitantes, não existe uma única solução
ótima, mas um conjunto de soluções não dominadas~--- a \emph{fronteira de
Pareto}~--- em que melhorar um objetivo implica piorar outro
\cite{deb2001multiobjective}. O NSGA-II \cite{deb2002nsga} é um algoritmo de
referência para esse cenário: combina ordenação rápida por não dominância, uma
medida de \emph{crowding distance} para preservar a diversidade ao longo da
fronteira e elitismo $(\mu+\lambda)$. A qualidade de uma fronteira é avaliada por
indicadores como o hipervolume (volume do espaço de objetivos dominado em relação
a um ponto de referência, que captura convergência e espalhamento em um único
número) e o espalhamento (\emph{spacing}, uniformidade da distribuição)
\cite{deb2001multiobjective,deb2002nsga}. O tratamento da preservação de
identidade como penalidade \emph{suave}, em vez de restrição rígida, alinha-se à
literatura de tratamento de restrições em algoritmos evolutivos
\cite{michalewicz1996constrained}.

\subsection{Balanceamento de jogos por busca}

A geração e o ajuste de conteúdo de jogos por busca (\emph{search-based procedural
content generation}) constituem um campo consolidado
\cite{togelius2011searchbased,yannakakis2018ai}. \cite{hom2007automatic}
automatizam o projeto de jogos de tabuleiro balanceados, avaliando configurações
por simulações de partidas. \cite{browne2010evolutionary}, no sistema Ludi,
evoluem regras de jogos completos guiados por uma bateria de critérios de
qualidade mensuráveis e validados por \emph{self-play}, culminando em jogos
efetivamente publicáveis~--- metodologia que inspira diretamente a etapa de
validação externa adotada neste trabalho. \cite{chen2014mmorpg} tratam o
balanceamento de classes de um \emph{MMORPG} por programação coevolutiva,
deixando as estratégias evoluírem contra adversários que também se adaptam, e
\cite{livingstone2006coevolution} explora coevolução em IA hierárquica para jogos
de estratégia. \cite{preuss2012balanced} buscam planos de jogo balanceados para um
cenário competitivo. A maior parte desses trabalhos otimiza o equilíbrio
diretamente; a contribuição específica deste trabalho é colocar o equilíbrio em
\emph{tensão explícita} com a preservação de identidade e mapear esse
\emph{trade-off}, em vez de tratá-lo como restrição implícita.

\subsection{Arquétipos e o ciclo de vantagens}
\label{sec:arquetipos}

A teoria de jogos competitivos descreve a saúde de um \emph{metagame} em termos de
relações de vantagem não transitivas~--- estruturas de pedra-papel-tesoura em que
cada opção vence algumas e perde para outras, o que impede a existência de uma
estratégia dominante
\cite{gingold2006rps,elias2012characteristics,sirlin2006playing}. Em jogos de
luta, esses papéis costumam ser descritos, no vocabulário corrente da comunidade
competitiva, por arquétipos reconhecíveis: \emph{rushdown} (agressor de curta
distância), \emph{zoner} (controlador de espaço), \emph{grappler} (especialista em
agarrões de alto dano), \emph{combo master} (conversor de acertos em sequências) e
\emph{turtle} (defensor que pune erros).

É importante delimitar o estatuto desses conceitos neste trabalho. A literatura
citada fundamenta a \emph{estrutura não transitiva} de forma geral, mas não define
um conjunto específico de arquétipos nem um ciclo de vantagens fechado. O
\emph{roster} de cinco personagens, os valores canônicos de seus atributos
(Tabela~\ref{tab:canonicos}) e o ciclo de vantagens correspondente são uma
\textbf{operacionalização construída pelos autores}, informada pela convenção da
comunidade. Eles cumprem três papéis: semente da população, \emph{baseline} de
medição de desvio e \emph{hipótese} de uma estrutura preservável. Não são tratados
como uma lei empírica a ser reproduzida pelo sistema. Essa delimitação reforça a
não circularidade: como o ciclo é uma hipótese dos próprios autores, verificar se
ele sobrevive sob o modelo é um achado sobre o modelo, e não a confirmação de um
fato externo. Cinco arquétipos bastam para um ciclo fechado, em que cada um vence
dois e perde para dois.

\section{Metodologia}

O sistema possui duas camadas independentes, orquestradas pelo algoritmo
evolutivo: uma \emph{camada de simulação} (combate 1\,vs\,1, tick a tick) e uma
\emph{camada evolutiva} (AG escalar e NSGA-II). A unidade de evolução é o
\emph{conjunto dos cinco personagens} (um por arquétipo), pois a taxa de vitória
de qualquer personagem depende simultaneamente dos outros quatro.

\subsection{Representação: arquétipos e cromossomo}

Cada indivíduo da população é um \emph{roster} completo de cinco personagens. Cada
personagem é descrito por 10 genes~--- 7 atributos numéricos (HP, dano,
\emph{cooldown} de ataque, alcance, velocidade, \emph{stun} e \emph{knockback}) e
3 pesos comportamentais ($w_{retreat}$, $w_{defend}$, $w_{aggressiveness}$)~---
totalizando 50 genes por indivíduo. Os pesos governam a política de combate
(Seção~\ref{sec:combate}). A Tabela~\ref{tab:canonicos} apresenta os valores
canônicos dos atributos, construídos conforme a Seção~\ref{sec:arquetipos}. Os
limites (\emph{bounds}) de cada gene foram calibrados de forma estreita~--- por
exemplo, $\mathrm{HP}\in [250,450]$ e $\mathrm{dano}\in [15,30]$~--- a fim de
eliminar modos de exploração degenerados pelo algoritmo, como personagens com
durabilidade arbitrariamente superior à do \emph{turtle}.

\begin{table}[ht]
\centering
\caption{Valores canônicos dos cinco arquétipos (7 atributos numéricos),
construídos pelos autores (Seção~\ref{sec:arquetipos}). Servem como semente da
população e \emph{baseline} de medição de desvio; não constituem restrições. O
\emph{stun} é expresso como \emph{fração} do \emph{cooldown} do próprio atacante
(Seção~\ref{sec:combate}), não como valor absoluto.}
\label{tab:canonicos}
\small
\begin{tabular}{lrrrrrrr}
\toprule
Arquétipo & HP & Dano & Cool. & Alc. & Vel. & Stun & Knock. \\
\midrule
Zoner        & 300 & 20 & 4 & 18 & 2{,}5 & 0{,}10 & 2{,}0 \\
Rushdown     & 320 & 16 & 1 & 10 & 5{,}0 & 0{,}10 & 1{,}0 \\
Combo Master & 350 & 18 & 3 & 10 & 3{,}0 & 0{,}55 & 0{,}5 \\
Grappler     & 400 & 27 & 4 &  8 & 2{,}0 & 0{,}30 & 0{,}5 \\
Turtle       & 450 & 15 & 5 & 13 & 1{,}5 & 0{,}20 & 1{,}0 \\
\bottomrule
\end{tabular}
\end{table}

\subsection{Modelo de combate}
\label{sec:combate}

O combate é uma simulação tick a tick, 1\,vs\,1, em um campo unidimensional de 100
unidades, com distância inicial de 50 e resolução em sub-ticks
($\texttt{TICK\_SCALE}=5$) para suavizar o espaço de \emph{fitness}. A cada
sub-tick, a ação de cada lutador~--- ATTACK, ADVANCE, RETREAT ou DEFEND~--- é
decidida em duas fases, \emph{intenção} seguida de \emph{execução} (um lutador sob
\emph{stun} perde a ação):

\begin{description}
\item[Fase 1 --- intenção.] Se está \emph{fora} do próprio alcance, a intenção é
ignorada e o lutador faz um ADVANCE incondicional (\emph{neutral game}, aproxima-se).
Se está \emph{em} alcance e não há intenção vigente, ele \emph{amostra} uma intenção
em $\{$FRENTE, RECUAR, GUARDA$\}$ com probabilidade proporcional aos pesos
$(w_{aggressiveness}, w_{retreat}, w_{defend})$ e a mantém pelos próximos
$\texttt{ACTION\_PERSISTENCE}=10$ sub-ticks (comprometimento/inércia estratégica).
\item[Fase 2 --- execução.] A intenção vigente é mapeada na ação concreta do
sub-tick: FRENTE $\rightarrow$ ATTACK se o \emph{cooldown} está pronto, senão
ADVANCE (pressiona sem desperdiçar o \emph{cooldown}); RECUAR $\rightarrow$ RETREAT
se ainda há espaço para recuar, senão DEFEND; GUARDA $\rightarrow$ DEFEND.
\end{description}

A amostragem da intenção é a \emph{única} fonte de estocasticidade do laço~--- não
há termo de hesitação nem tratamento especial de encurralamento. Como os pesos
agem de forma contínua, uma variação infinitesimal em qualquer peso produz
variação proporcional na probabilidade da intenção, o que fornece ao algoritmo um
gradiente contínuo sobre esses genes; uma comparação categórica (só a ordem dos
pesos importando) tornaria as magnitudes invisíveis à seleção.

A resolução de dano é determinística e \emph{flat}: o dano é o próprio atributo do
atacante, sem redução passiva; o único modificador é a ação DEFEND do alvo, que
reduz o dano recebido em 40\%. O \emph{stun} efetivo é
$\mathrm{round}\!\big(\textit{stun}\times\mathrm{round}(\textit{cooldown}\times5)\big)$,
em que o gene $\textit{stun}\in[0;\,0{,}6]$ é uma \emph{fração} do \emph{cooldown}
do próprio atacante (em sub-ticks). Como esse limite é inferior a 1, o \emph{stun}
aplicado é sempre estritamente menor que o \emph{cooldown}, o que garante sempre
uma janela livre entre golpes e dispensa um teto explícito. Uma luta termina por
nocaute ou, ao atingir o limite de duração, pela maior fração de HP remanescente.

\subsection{Definição e medição de identidade}
\label{sec:identidade}

Um arquétipo não é definido por valores numéricos exatos, mas por
\emph{invariantes qualitativos}: o \emph{zoner} é aquele que controla o espaço à
distância; o \emph{turtle}, aquele que absorve dano e pune erros. Para tornar a
preservação de identidade verificável, este trabalho a instrumenta por quatro
lentes complementares, três das quais são puramente de medição e uma das quais
(o desvio) também integra a função objetivo (Seção~\ref{sec:fitness}):

\begin{description}
\item[Estrutural e comportamental (validador).] Um conjunto de \textbf{21 asserções
de ranqueamento}, em três camadas, verifica se cada personagem ocupa, dentro do
\emph{roster}, a posição relativa que define seu arquétipo. As duas primeiras
camadas são \emph{estruturais} (17 asserções sobre os genes): 12 inter-personagem
(por exemplo, o \emph{rushdown} deve ter a maior velocidade e o menor
\emph{cooldown}; o \emph{turtle}, o maior HP, o maior \emph{cooldown} e a maior
propensão a defender; o \emph{grappler}, o maior dano) e 5 intra-personagem sobre
ordenações internas normalizadas (por exemplo, no \emph{zoner}, alcance normalizado
maior que velocidade normalizada). A terceira camada é \emph{comportamental} (4
asserções): afere se a identidade se traduz em jogo efetivo~--- o \emph{zoner} luta
da maior distância média, o \emph{rushdown} conecta o maior número de ataques, o
\emph{turtle} é quem mais \emph{escolhe} defender e o \emph{combo master} é quem
mais aplica \emph{stun}. O \emph{grappler} não recebe asserção comportamental: como
o modelo não representa a mecânica de \emph{grab}, sua identidade não tem expressão
comportamental distinta~--- um achado honesto de que 4 das 5 identidades se
manifestam funcionalmente. Por se basear em ranqueamento, e não em distância a
valores específicos, o instrumento é \emph{invariante à magnitude exata} dos genes:
afere a identidade no sentido qualitativo que a define.

\item[Comportamental (\emph{fingerprint}).] Para cada personagem, agregado sobre
seus quatro confrontos, registra-se o perfil de jogo: a distribuição de ações
(ATTACK/ADVANCE/RETREAT/DEFEND), a fração de tempo fora de alcance (espaçamento) e
a fração de tempo sob \emph{stun}. A comparação entre o perfil canônico e o
evoluído (em pontos percentuais) verifica se o personagem evoluído \emph{ainda
joga} como o arquétipo~--- se o \emph{zoner} continua a recuar e espaçar, se o
\emph{turtle} continua a defender.

\item[Numérico (\emph{drift}).] A distância euclidiana normalizada de cada
personagem ao seu perfil canônico (Seção~\ref{sec:fitness}). É a tradução contínua
de ``preservação de identidade'' e o único dos quatro instrumentos que também é
otimizado.

\item[Diferenciação.] A distância média par a par entre os cinco personagens,
expressa como razão em relação à do \emph{roster} canônico. Razão menor que 1
indica homogeneização (os personagens convergiram entre si); maior que 1 indica
que o elenco evoluído é mais distinto que o de partida.
\end{description}

Apenas o desvio entra na otimização; validador, \emph{fingerprint} e diferenciação
são exclusivamente de medição. Essa separação é o cerne da não circularidade: a
homogeneização e a preservação do ciclo são \emph{medidas}, nunca recompensadas
pela função objetivo.

\subsection{Função de avaliação: os dois objetivos}
\label{sec:fitness}

A avaliação de um indivíduo é feita por \emph{round-robin} completo: os
$\binom{5}{2}=10$ confrontos são simulados $\texttt{SIMS\_PER\_MATCHUP}=150$ vezes
cada. Dela derivam os dois objetivos do problema, que são exatamente os dois eixos
da pergunta de pesquisa.

\paragraph{Desvio (\emph{drift}).} Operacionaliza a preservação de identidade
como a distância euclidiana normalizada de cada personagem ao seu perfil
canônico, sobre os 10 genes, agregada pela média entre os cinco:
\begin{equation}
\textit{drift} = \frac{1}{5}\sum_{i=1}^{5}
\sqrt{\frac{1}{10}\!\left(\sum_{a}\!\Big(\tfrac{g_{a}-c_{a}}{g_{a}^{\max}}\Big)^{\!2}
+ \sum_{w}(g_{w}-c_{w})^{2}\right)} .
\end{equation}
É também o verdadeiro mecanismo anti-homogeneização da otimização, pois atrai cada
personagem para um perfil canônico distinto.

\paragraph{Desbalanço (\emph{dominance}).} Combina três sinais cegos à direção,
cada um agregado como raiz quadrática média (RMS), de modo que casos extremos pesem
mais que moderados. O termo \emph{primário} é o desbalanço \emph{global} por
personagem: para cada um dos cinco, $|\textit{WR}_{\textit{glob}}-0{,}5|/0{,}5$, em
que $\textit{WR}_{\textit{glob}}$ é a taxa de vitória (\emph{win rate}, WR) agregada
sobre os quatro oponentes. Seu ótimo é ``nenhum arquétipo domina o \emph{roster}''
e, crucialmente, ele \emph{não} força cada par a 50\%: um personagem a 50\% global
pode vencer dois confrontos e perder dois~--- exatamente o espaço em que o ciclo de
vantagens pode existir. Dois termos \emph{secundários}, ambos por par, regularizam a
qualidade: (i) um \emph{teto de hard-counter}, que penaliza apenas o excesso de
$|\textit{WR}_{\textit{par}}-0{,}5|$ acima de $\texttt{MATCHUP\_WR\_CAP}=0{,}15$,
mantendo as arestas do ciclo como \emph{vantagens} dentro de $[0{,}35;\,0{,}65]$ e
barrando counters esmagadores (como 100\%$\times$0\%); e (ii) a \emph{decisividade}
por luta fora de uma banda saudável $[0{,}10;\,0{,}20]$ (em que o vencedor encerra
com 20\% a 40\% de HP de folga). Sejam
\begin{equation}
g=\sqrt{\tfrac{1}{5}\textstyle\sum_{i=1}^{5}\!\Big(\tfrac{|\textit{WR}_{\textit{glob},i}-0{,}5|}{0{,}5}\Big)^{\!2}},\ \
c=\sqrt{\tfrac{1}{10}\textstyle\sum_{ij}\!\Big(\tfrac{\max(0,\,|\textit{WR}_{ij}-0{,}5|-0{,}15)}{0{,}35}\Big)^{\!2}},\ \
d=\sqrt{\tfrac{1}{10}\textstyle\sum_{ij}\textit{decis\_excess}_{ij}^{2}},
\end{equation}
\begin{equation}
\textit{dominance} = w_{g}\,g + w_{c}\,c + w_{d}\,d,\qquad
w_{g}=1{,}0,\quad w_{c}=w_{d}=0{,}5 .
\end{equation}
Os três sinais usam $|\textit{WR}-0{,}5|$ e $|\textit{score}-0{,}5|$~--- são
\emph{cegos à direção}: a penalidade não codifica \emph{qual} arquétipo deveria
vencer cada confronto, o que preserva a não circularidade.

Os três termos são complementares porque cada um, isolado, é insuficiente. O
balanço global é cego ao \emph{formato} de cada par: admite tanto um counter
esmagador (100\%$\times$0\%, com WR global ainda equilibrada) quanto um
\emph{blowout-coinflip}~--- situação em que cada lado vence de forma esmagadora em
metade das partidas, produzindo WR próxima de 50\% embora nenhuma luta seja de fato
equilibrada. O teto de \emph{hard-counter} barra o primeiro; a decisividade por
luta, medida como margem contínua de HP, captura o segundo. O ruído introduzido
pela amostragem de intenção torna a WR um gradiente utilizável: quando os
confrontos se aproximam do equilíbrio, pequenas variações de genes deslocam
continuamente as taxas de vitória.

\subsection{AG escalar e NSGA-II}

O \textbf{AG escalar} otimiza a soma ponderada
$\textit{fitness} = -(\lambda_{d}\cdot\textit{drift} +
\lambda_{D}\cdot\textit{dominance})$, com $\lambda_{d}=\lambda_{D}=1{,}0$, de modo
que identidade e equilíbrio pesem na mesma escala. Ele entrega \emph{um} ponto do
\emph{trade-off}, correspondente à abordagem convencional de soma ponderada. O
\textbf{NSGA-II} \cite{deb2002nsga} dispensa os pesos e otimiza o par
$(\textit{dominance}, \textit{drift})$ diretamente, mapeando \emph{toda} a
fronteira~--- de ``preserva e desequilibra'' a ``equilibra e homogeneíza''. É o
artefato central deste trabalho, por tornar o \emph{trade-off} explícito e
visualizável.

Ambos compartilham os operadores: seleção por torneio (tamanho 3),
\emph{crossover} por bloco de personagem (cada personagem do filho é herdado
integralmente de um dos pais, o que preserva a coerência interna do arquétipo),
mutação gaussiana com desvio quatro vezes menor nos pesos do que nos atributos
(inércia estratégica) e elitismo de 10\%. A população é de 300 indivíduos,
evoluídos por até 150 gerações. Da fronteira final, quatro representantes são
extraídos: \texttt{best\_dominance} (maior equilíbrio), \texttt{best\_drift}
(maior fidelidade ao canônico), \texttt{knee\_point} (melhor compromisso, ponto de
máxima curvatura) e \texttt{ideal\_point} (mais próximo da utopia).

\subsection{Protocolo experimental e reprodutibilidade}
\label{sec:protocolo}

Como um algoritmo evolutivo é estocástico, uma única execução constitui uma
amostra, e não um resultado \cite{eiben2015introduction,deb2001multiobjective}. O
protocolo experimental, incorporado da literatura, prevê três instrumentos: (i)
\emph{múltiplas execuções independentes} com estatística agregada~--- média $\pm$
desvio-padrão, taxa de sucesso por confronto e fração de sementes que equilibram
todos os 10 confrontos; (ii) \emph{métricas de qualidade da fronteira}~---
hipervolume com ponto de referência fixo $(2{,}0,\,1{,}0)$ e espalhamento; e (iii)
\emph{validação externa ao} \emph{fitness}, ao estilo Ludi
\cite{browne2010evolutionary}, que reavalia um indivíduo fixo sob sementes de
avaliação inteiramente novas, fora do laço de otimização, e emite um veredito de
robustez por confronto.

A reprodutibilidade exige cuidado específico: o combate amostra com o gerador
interno do compilador \emph{just-in-time} (Numba), independente do gerador do
Python. A semeadura é, por isso, feita por rotina dedicada que reinicializa o
estado ao mesmo \emph{seed} base em toda avaliação (\emph{Common Random Numbers}),
de modo que as diferenças de \emph{fitness} reflitam os genes, e não o sorteio.
Salvo indicação contrária, as execuções individuais reportadas a seguir usam a
semente \ph{seed}; a agregação multissemente usa \ph{N} sementes independentes.

\section{Resultados}

Esta seção segue o fio condutor: (i) estabelecer o \emph{baseline} canônico, tanto
em identidade quanto em equilíbrio; (ii) mostrar que a otimização direta do
equilíbrio o atinge, mas ao custo de homogeneização, o que motiva o tratamento
bi-objetivo; (iii) apresentar a fronteira de Pareto como o mapa do
\emph{trade-off}; e (iv) consolidar a evidência com estatística agregada e
validação de robustez.

\subsection{O \emph{baseline} canônico: identidade validada, equilíbrio ausente}

Antes da evolução, o \emph{roster} canônico possui identidade perfeita por
construção. O validador confirma esse ponto de forma objetiva:
\textbf{o roster canônico satisfaz as 21 de 21 asserções} de identidade (17
estruturais, por construção, e 4 comportamentais), e o desvio é nulo
($\textit{drift}=0$). O \emph{fingerprint} comportamental é coerente
com os conceitos~--- o \emph{rushdown} avança em \ph{fp-rush} dos ticks, o
\emph{turtle} defende em \ph{fp-turtle} e o \emph{zoner} recua em \ph{fp-zoner}~---,
o que indica que os pesos produzem identidade comportamental visível, e não ruído
arbitrário.

Apesar dessa integridade de identidade, o roster é marcadamente desequilibrado. A
validação externa atribui-lhe $\textit{dominance}=\ph{dom-can-ext}$ e
\textbf{\ph{eq-can}} confrontos equilibrados: o Grappler vence \ph{wr-grap} de suas
lutas, enquanto o Zoner vence \ph{wr-zon}; \ph{blow-can} dos dez confrontos são
unilaterais (\emph{blowouts}). Esse resultado é, em si, um achado: o ciclo de vantagens canônico não é
trivialmente preservado em um modelo de combate determinístico~--- apenas
\ph{cycle-can} dos dez pares do ciclo se sustentam~---, o que sugere que a
estrutura não transitiva dos jogos de luta reais depende parcialmente de mecânicas
(combos encadeados, variância) deliberadamente ausentes do modelo. O
\emph{baseline} define o ponto de partida contra o qual todo indivíduo evoluído é
lido.

\subsection{A otimização direta do equilíbrio custa identidade}

O AG escalar, com pesos iguais, persegue um único ponto e o empurra para o
equilíbrio. Ele alcança $\textit{dominance}=\ph{dom-ag}$ e \textbf{\ph{eq-ag}}
confrontos equilibrados, com WR em $[40\%,60\%]$. Esse equilíbrio quase total tem,
porém, um custo de identidade visível e mensurável pelos instrumentos
\emph{post hoc}: o desvio sobe para \ph{drift-ag} e, sobretudo, a diferenciação cai
para uma razão de \ph{diff-ag}~--- sinal de homogeneização, isto é, os cinco
personagens convergiram parcialmente entre si. O resultado mostra empiricamente
que igualar as taxas de vitória, por si só, tende a dissolver a distinção entre
arquétipos~--- precisamente a tensão que a pergunta de pesquisa identifica~--- e
motiva tratar identidade e equilíbrio como objetivos separados, em vez de fundi-los
em uma única medida escalar.

\subsection{A fronteira de Pareto torna o \emph{trade-off} explícito}

A Figura~\ref{fig:pareto} apresenta a fronteira de Pareto obtida pelo NSGA-II após
150 gerações: \ph{n-front} soluções não dominadas, formando uma curva bem
distribuída (espalhamento de \ph{spacing}), com hipervolume de \ph{hv} em relação
ao ponto de referência $(2{,}0,\,1{,}0)$. A fronteira responde diretamente à
pergunta de pesquisa: evidencia que equilíbrio e identidade \emph{são} objetivos
conflitantes neste domínio e quantifica \emph{quanto} de um se troca pelo outro. O
extremo de identidade preservada (\texttt{best\_drift}) coincide com o canônico
($\textit{drift}=0$, $\textit{dominance}=\ph{dom-can}$); à medida que se admite
desvio, a penalidade de desbalanço diminui de forma monotônica, até o extremo de
equilíbrio (\texttt{best\_dominance}: $\textit{dominance}=\ph{dom-bd}$,
$\textit{drift}=\ph{drift-bd}$). O \emph{knee point}, melhor compromisso, situa-se
em $(\ph{knee-dom},\,\ph{knee-drift})$.

\begin{figure}[ht]
\centering
\includegraphics[width=.8\textwidth]{pareto_front.png}
\caption{Fronteira de Pareto (desbalanço $\times$ desvio) obtida pelo NSGA-II.
Cada ponto é um \emph{roster} não dominado. Os representantes destacados marcam os
extremos (\texttt{best\_dominance}, \texttt{best\_drift}), o \emph{knee point}
(melhor compromisso) e o \emph{ideal point}. A curva torna explícita a troca entre
equilíbrio e preservação de identidade.}
\label{fig:pareto}
\end{figure}

\subsection{Equilíbrio sem homogeneização: a leitura dos extremos}

A Tabela~\ref{tab:tradeoff} resume três pontos do \emph{trade-off}, medidos pelos
quatro instrumentos de identidade. A comparação é o ponto central do trabalho. No
extremo de identidade perfeita (canônico), o equilíbrio é nulo (\ph{eq-can}). No
ponto do AG escalar, o equilíbrio é quase total (\ph{eq-ag}), mas a diferenciação
cai (\ph{diff-ag}) e o validador registra apenas \ph{val-ag} de 21
asserções preservadas. O representante \texttt{best\_dominance} do NSGA-II ocupa a posição
mais relevante: pagando um desvio modesto de \ph{drift-bd}, reduz o desbalanço de
\ph{dom-can} para \ph{dom-bd} e atinge \textbf{\ph{eq-bd}} confrontos
equilibrados~--- e, decisivamente, a diferenciação par a par \emph{aumenta} em
relação ao canônico (razão \ph{diff-bd}), com \ph{val-bd} de 21 asserções
de identidade preservadas. Trata-se de equilíbrio substancial \emph{sem}
homogeneização. Mesmo sem ser codificado na função objetivo, \ph{cycle-bd} dos dez
pares do ciclo de vantagens canônico sobrevivem nesse indivíduo, como métrica
\emph{post hoc}.

\begin{table}[ht]
\centering
\caption{Três pontos do \emph{trade-off} equilíbrio $\times$ identidade, medidos
por métricas \emph{post hoc}. As setas indicam a direção de melhoria. ``Conf.
eq.''~= pares com WR em $[40\%,60\%]$; ``Difer.''~= razão da diferenciação par a
par em relação à do canônico ($>1$ indica maior distinção); ``Valid.''~= asserções
de identidade preservadas (de 21; 17 estruturais + 4 comportamentais); ``Ciclo''~=
pares do ciclo canônico mantidos \emph{post hoc}.}
\label{tab:tradeoff}
\small
\begin{tabular}{lcccccc}
\toprule
Configuração & \emph{dominance} $\downarrow$ & \emph{drift} $\downarrow$ &
Conf. eq. & Difer. & Valid. & Ciclo \\
\midrule
Canônico (\texttt{best\_drift})    & \ph{dom-can} & 0{,}00 & \ph{eq-can} & 1{,}00 & 21/21 & \ph{cycle-can} \\
AG escalar ($\lambda_d{=}\lambda_D$) & \ph{dom-ag} & \ph{drift-ag} & \ph{eq-ag} & \ph{diff-ag} & \ph{val-ag} & \ph{cycle-ag} \\
NSGA-II \texttt{best\_dominance}   & \ph{dom-bd} & \ph{drift-bd} & \ph{eq-bd} & \ph{diff-bd} & \ph{val-bd} & \ph{cycle-bd} \\
\bottomrule
\end{tabular}
\end{table}

A lição é que empurrar o equilíbrio até o limite custa identidade, ao passo que um
ponto intermediário da fronteira captura a maior parte do ganho de equilíbrio
preservando~--- ou até ampliando~--- a distinção entre arquétipos. A fronteira de
Pareto é justamente o que permite \emph{escolher} conscientemente esse ponto, em
vez de aceitar o resultado fixo de uma soma ponderada.

\subsection{Estatística agregada e robustez}

Para distinguir o que é estrutural do que é amostral, os resultados acima são
consolidados sobre \ph{N} execuções independentes. Em média, o NSGA-II aterrissa
em $\textit{dominance}=\ph{dom-mean}\pm\ph{dom-std}$ e
$\textit{drift}=\ph{drift-mean}\pm\ph{drift-std}$, com hipervolume médio de
\ph{hv-mean}$\pm$\ph{hv-std}. A taxa de sucesso por confronto identifica quais
pares o sistema equilibra de forma consistente: \ph{sr-resumo}. Em
\textbf{\ph{frac-balance}} das execuções, o representante de maior equilíbrio
equilibra todos os 10 confrontos.

A robustez do indivíduo \texttt{best\_dominance} é então verificada fora do laço de
otimização, sob \ph{ext-N} sementes de avaliação inteiramente novas
(Seção~\ref{sec:protocolo}). O veredito é que \ph{rob-bd} dos 10 confrontos
permanecem equilibrados em todas as condições, o que indica que o equilíbrio
reportado \ph{rob-verdito} um artefato de sobreajuste ao \emph{fitness}.

\subsection{Síntese dos achados}

Os resultados sustentam uma resposta afirmativa, ainda que matizada, à pergunta de
pesquisa: o equilíbrio competitivo \emph{pode} ser alcançado sem destruir as
identidades funcionais, mas há um \emph{trade-off} mensurável, e o grau de
equilíbrio buscado determina quanto de identidade se preserva. O equilíbrio
substancial (\ph{eq-bd}) é compatível com identidade preservada~--- e até
reforçada, conforme a diferenciação~--- ao passo que o equilíbrio total
(\ph{eq-ag}), nesta configuração, exige alguma homogeneização.

\section{Limitações}

O alcance dos resultados deve ser lido à luz de quatro delimitações. Primeira, o
modelo de combate é uma simplificação dos jogos de luta reais: não modela
\emph{frames} de inicialização e recuperação por golpe, \emph{mix-ups},
\emph{neutral game} nem oclusão. O próprio achado de que o ciclo canônico não
sobrevive ao modelo determinístico é evidência de que essas mecânicas omitidas têm
papel estrutural; por isso, o objeto de estudo não é a fidelidade a um jogo
específico, mas o \emph{mecanismo do} \emph{trade-off} entre equilíbrio e
identidade. Segunda, o estudo emprega cinco arquétipos, enquanto jogos
reais costumam ter dez ou mais. Terceira, o \emph{round-robin} pressupõe que todos
os arquétipos sejam jogados em igual proporção, e não modela o \emph{matchmaking},
em que jogadores escolhem confrontos favoráveis. Quarta, o \emph{trade-off} é
formulado como soma ponderada (escalar) ou Pareto bidimensional (NSGA-II); outras
formulações são possíveis e ficam como trabalho futuro.

\section{Conclusão e Trabalhos Futuros}
\label{sec:conclusao}

Este artigo formulou o balanceamento de personagens em jogos competitivos como um
problema bi-objetivo~--- equilíbrio \emph{versus} preservação de identidade~--- e o
investigou com um AG escalar e com o NSGA-II, sob a premissa metodológica de
\emph{medir}, e não \emph{impor}, a identidade, o que evita uma pergunta de
pesquisa circular. Os resultados mostram que o \emph{roster} canônico é fortemente
desequilibrado, que a otimização direta do equilíbrio o homogeneíza, e que a
fronteira de Pareto torna o \emph{trade-off} explícito e mensurável, permitindo
obter equilíbrio substancial enquanto se \emph{aumenta} a diferenciação entre
arquétipos. A separação entre objetivos otimizados e métricas \emph{post hoc}
(ciclo, diferenciação, comportamento e estrutura) é o que preserva a validade
científica da resposta.

Como trabalhos futuros, destacam-se três direções. A primeira é o \emph{teste de
estresse} do equilíbrio por coevolução
\cite{chen2014mmorpg,livingstone2006coevolution}, para verificar se o equilíbrio
sobrevive a um oponente que também se adapta. A segunda é a reformulação por
\emph{quality-diversity} (MAP-Elites e \emph{novelty search})
\cite{mouret2015mapelites,lehman2011novelty}, que substituiria a busca por uma
única solução pela iluminação de \emph{quantas} configurações distintas de
\emph{roster} atingem equilíbrio, atacando ainda mais diretamente a tese de
equilíbrio \emph{sem} perda de identidade. A terceira é a análise causal de
mecânicas por \emph{restricted play} \cite{hom2007automatic}, que permitiria
diagnosticar quais mecânicas sustentam ou degeneram cada confronto, com vistas a
reaproximar o modelo das estruturas observadas em jogos reais.

\bibliographystyle{sbc}
\bibliography{referencias}

\end{document}
